% ═══════════════════════════════════════════════════════════════════════
\section{Compiler et corriger les erreurs}
% ═══════════════════════════════════════════════════════════════════════

\begin{frame}{Ce qui se passe quand on compile}
  \small
  \begin{itemize}
    \item Le compilateur (pdfLaTeX sur Overleaf) lit le fichier
      \texttt{.tex} de haut en bas, exécute les commandes et produit un
      PDF.
    \item Il écrit aussi un journal (le « log ») : avertissements et
      erreurs, avec le numéro de la ligne en cause.
    \item Avertissement (jaune) : le PDF est produit, mais quelque chose
      est à vérifier (une ligne trop longue, une référence inconnue).
    \item Erreur (rouge) : la compilation s'est arrêtée ou a dû deviner.
      Il faut la corriger, même si un PDF s'affiche.
  \end{itemize}
  \vspace{4pt}
  \begin{Rem}{Sur Overleaf}
    \small Le bouton à côté de Recompiler (Journaux et fichiers produits,
    ou Logs and output files) liste les problèmes ; un clic sur une erreur
    amène à la ligne correspondante dans l'éditeur.
  \end{Rem}
\end{frame}

\begin{frame}{Les quatre erreurs que vous verrez le plus}
  \footnotesize
  \begin{center}
  \begin{tabular}{p{3.6cm}p{4.2cm}p{4.9cm}}
    \toprule
    Message & Cause probable & Correction \\
    \midrule
    \texttt{Undefined control sequence} & une commande mal écrite
      (\cmd{fracc}) ou dont l'extension n'est pas chargée & corriger
      l'orthographe, ou ajouter le \cmd{usepackage} \\
    \texttt{Missing \$ inserted} & des mathématiques hors du mode
      mathématique (\texttt{x\^{}2} dans le texte, un \texttt{\_}) &
      entourer de \texttt{\$...\$} \\
    \texttt{Runaway argument}, \texttt{Missing \} inserted} & une
      accolade oubliée & compter les \texttt{\{} et les \texttt{\}} \\
    \texttt{\textbackslash begin\{x\} ended by \textbackslash end\{y\}} &
      environnements mal fermés ou croisés & refermer dans l'ordre
      inverse d'ouverture \\
    \texttt{File `xyz.sty' not found} & extension inexistante (souvent
      inventée par l'IA) & retirer la ligne ou demander une
      alternative \\
    \bottomrule
  \end{tabular}
  \end{center}
  \small Le numéro de ligne s'écrit \texttt{l.12} : c'est là que \LaTeX\
  s'est aperçu du problème --- la cause est à cet endroit ou juste avant.
\end{frame}

\begin{frame}{Une méthode pour ne pas paniquer}
  \small
  \begin{enumerate}
    \item Lire la première erreur seulement, et sa ligne.
    \item Regarder cette ligne et celle d'avant : accolade, \texttt{\$},
      orthographe, \cmd{end} manquant ?
    \item Si on ne trouve pas : mettre des \texttt{\%} devant le bloc
      suspect pour le neutraliser, recompiler, réduire le bloc.
    \item Toujours pas ? Demander à l'IA en lui donnant le message
      d'erreur complet et les lignes concernées (prompt plus loin).
  \end{enumerate}
  \vspace{4pt}
  \begin{AuTableau}[Le bon réflexe]
    \small Compiler après chaque petite modification. Une erreur dans
    trois lignes que l'on vient d'écrire se trouve en dix secondes ; dans
    trois pages, elle peut prendre une demi-heure.
  \end{AuTableau}
\end{frame}

\begin{frame}[fragile]{À vous}
  \begin{columns}[c]
    \begin{column}{0.52\textwidth}
      \begin{Exercice}{Réparer}
        \small Copiez ces lignes dans le corps de votre document,
        compilez, lisez les messages, et corrigez les cinq erreurs.
      \end{Exercice}
    \end{column}
    \begin{column}{0.45\textwidth}
\begin{CodeTeX}[title={Cinq erreurs cachées}]
\section{Exercice
Calculer x^2 si $x = 3$.
\begin{itemize}
  \item La fraction $\frac{1}{2$.
  \item Le rabais : 20%.
\end{enumerate}
\end{CodeTeX}
    \end{column}
  \end{columns}
\end{frame}
